\documentclass[12pt]{article}
\usepackage[margin=1in]{geometry}
\usepackage{amsmath}
\usepackage{amssymb}
\usepackage{booktabs}
\usepackage{hyperref}
\usepackage{listings}
\usepackage{xcolor}
\usepackage{titlesec}
\usepackage{parskip}
\usepackage{float}
\usepackage{enumitem}
\usepackage{tcolorbox}

\tcbuselibrary{skins, breakable}

\definecolor{codebg}{RGB}{245,245,245}
\definecolor{notebg}{RGB}{255,248,220}
\definecolor{successbg}{RGB}{235,255,235}
\definecolor{warnbg}{RGB}{255,235,235}

\hypersetup{
    colorlinks=true,
    linkcolor=blue,
    urlcolor=blue
}

\lstset{
    basicstyle=\ttfamily\small,
    backgroundcolor=\color{codebg},
    frame=single,
    breaklines=true,
    keywordstyle=\color{blue},
    commentstyle=\color{green!60!black},
    stringstyle=\color{orange},
    numbers=left,
    numberstyle=\tiny\color{gray},
    numbersep=5pt
}

\newtcolorbox{notebox}{
    colback=notebg,
    colframe=orange!70!black,
    title=\textbf{Note},
    fonttitle=\small
}

\newtcolorbox{successbox}{
    colback=successbg,
    colframe=green!60!black,
    title=\textbf{Result},
    fonttitle=\small
}

\newtcolorbox{warnbox}{
    colback=warnbg,
    colframe=red!70!black,
    title=\textbf{Challenge},
    fonttitle=\small
}

\title{\textbf{Phase 2 Simulation Report} \\
\large CBF Safety Filter and Hybrid State Machine \\
\normalsize Surface-Proximate Drone Docking Project}
\author{Sean Dai}
\date{March 14, 2026}

\begin{document}

\maketitle
\tableofcontents
\newpage

% -------------------------------------------------------
\section{Overview}
% -------------------------------------------------------

Phase 2 added a mathematically rigorous safety layer on top of the
Phase 1 simulation. The nominal cascaded PD controller from Phase 1
continues to run unchanged. A new Control Barrier Function (CBF)
safety filter intercepts the nominal control input and modifies it
only when necessary to prevent unsafe states. A hybrid state machine
governs which operating mode the drone is in based on the Dockability
Margin $M$ calibrated in Phase 1.

\textbf{Phase 2 produced three verified results:}

\begin{enumerate}
    \item The CBF filter activates at the correct $M$ threshold and
          modifies control inputs to prevent wall collision.
    \item The state machine transitions correctly through all four
          modes during a wall approach.
    \item The CBF-filtered trajectory handles abort and retreat cycles
          safely while the unfiltered trajectory remains dangerously
          close to the wall without recovery.
\end{enumerate}

% -------------------------------------------------------
\section{New Dependencies and File Structure}
% -------------------------------------------------------

\subsection{New Library}

Phase 2 required one new Python library for the quadratic program solver:

\begin{lstlisting}[language=bash, numbers=none]
pip install cvxpy==1.8.1
\end{lstlisting}

\subsection{Updated File Structure}

Phase 2 added two new files and modified four existing ones.

\begin{center}
\begin{tabular}{lll}
\toprule
\textbf{File} & \textbf{Status} & \textbf{Responsibility} \\
\midrule
\texttt{params.py}         & Modified & CBF gains and retreat parameters added \\
\texttt{dynamics.py}       & Modified & Wall clamp added to prevent wall crossing \\
\texttt{logger.py}         & Modified & Mode, CBF flag, and T\_nom fields added \\
\texttt{visualize.py}      & Modified & Mode transition plot added \\
\texttt{run\_sim.py}       & Modified & CBF and state machine integrated \\
\texttt{cbf.py}            & New      & CBF safety filter \\
\texttt{state\_machine.py} & New      & Hybrid docking state machine \\
\bottomrule
\end{tabular}
\end{center}

% -------------------------------------------------------
\section{New Parameters}
% -------------------------------------------------------

The following parameters were added to \texttt{params.py}:

\begin{center}
\begin{tabular}{lll}
\toprule
\textbf{Parameter} & \textbf{Value} & \textbf{Description} \\
\midrule
\texttt{cbf\_gamma\_d}  & 0.5  & Class-$\mathcal{K}$ gain for wall clearance \\
\texttt{cbf\_gamma\_v}  & 0.5  & Class-$\mathcal{K}$ gain for velocity \\
\texttt{cbf\_gamma\_th} & 1.0  & Class-$\mathcal{K}$ gain for attitude \\
\texttt{cbf\_gamma\_u}  & 1.0  & Class-$\mathcal{K}$ gain for actuator \\
\texttt{x\_retreat}     & 1.5  & Retreat target position (m) \\
\texttt{retreat\_speed} & 0.05 & Retreat approach speed (m/s) \\
\bottomrule
\end{tabular}
\end{center}

% -------------------------------------------------------
\section{State Machine Design}
% -------------------------------------------------------

\subsection{Four Operating Modes}

The state machine reads $M$ at every timestep and returns one of four
mode strings. Each mode determines whether the CBF filter is active
and what target position the controller tracks.

\begin{center}
\begin{tabular}{llll}
\toprule
\textbf{Mode} & \textbf{Entry condition} & \textbf{CBF} &
\textbf{Target} \\
\midrule
\texttt{free}    & $M \geq M_\mathrm{safe}$   & Off &
    Approach lambda \\
\texttt{surface} & $M < M_\mathrm{safe}$      & On  &
    Approach lambda \\
\texttt{docking} & $M < M_\mathrm{dock}$      & On  &
    Approach lambda \\
\texttt{abort}   & $M < M_\mathrm{abort}$     & On  &
    $x_\mathrm{retreat}$ \\
\bottomrule
\end{tabular}
\end{center}

\subsection{Hysteresis}

A hysteresis band of 0.05 was implemented to prevent rapid
mode-switching at threshold boundaries. For example, the drone does
not re-enter \texttt{free} mode until $M$ rises to
$M_\mathrm{safe} + 0.05 = 0.65$, not just above 0.6. This prevents
chattering behavior at threshold crossings.

\subsection{Abort Recovery}

Once abort is triggered, the drone targets $x_\mathrm{retreat}$
permanently until $M$ recovers above $M_\mathrm{safe} + 0.1 = 0.7$.
At that point \texttt{abort\_triggered} is reset and the approach
lambda resumes from the current position. This allows multiple
approach-abort-retreat cycles within a single simulation run.

\subsection{Verification Test}

\begin{lstlisting}[language=Python]
from params import PARAMS
from state_machine import StateMachine

sm = StateMachine(PARAMS)
state = [1.5, 1.5, 0.0, 0.0, 0.0, 0.0]

print(sm.update(0.8, state, PARAMS))   # free
print(sm.update(0.5, state, PARAMS))   # surface
print(sm.update(0.3, state, PARAMS))   # docking
print(sm.update(0.1, state, PARAMS))   # abort
\end{lstlisting}

\textbf{Result:} All four modes returned correctly in sequence.

% -------------------------------------------------------
\section{CBF Safety Filter Design}
% -------------------------------------------------------

\subsection{Architecture Decision}

The initial CBF design used a \texttt{cvxpy} quadratic program with
thrust and torque as decision variables. This approach failed because
at $\theta = 0$ (level hover), $\sin\theta = 0$, making all
wall-related constraints trivially satisfied regardless of closing
velocity. The QP had no way to enforce safety through thrust alone
when the drone is flying level.

The final implementation instead operates at the outer loop level by
limiting the desired pitch angle based on wall distance and closing
velocity. Since pitch angle directly controls horizontal acceleration,
this is the correct level at which to enforce wall safety.

\subsection{Barrier Functions}

Two barrier functions are enforced:

\textbf{Barrier 1 — Wall clearance:}
\begin{equation}
    h_d(x) = d - d_{\min} \geq 0
\end{equation}
\begin{equation}
    \dot{h}_d = v_n \geq -\gamma_d(d - d_{\min})
\end{equation}

\textbf{Barrier 2 — Approach velocity:}
\begin{equation}
    h_v(x) = v_{n,\max} + v_n \geq 0
\end{equation}
\begin{equation}
    v_n \geq -v_{n,\max}
\end{equation}

The combined floor on closing velocity is:
\begin{equation}
    v_{n,\mathrm{floor}} = \max\left(-\gamma_d(d - d_{\min}),\;
    -v_{n,\max}\right)
\end{equation}

\subsection{Control Law}

If the current $v_n < v_{n,\mathrm{floor}}$, the CBF computes a
required deceleration and converts it to a safe pitch angle:

\begin{equation}
    a_{x,\mathrm{req}} = \gamma_d(v_{n,\mathrm{floor}} - v_n)
\end{equation}

\begin{equation}
    \theta_\mathrm{safe} = \arcsin\!\left(
    \mathrm{clip}\!\left(\frac{-a_{x,\mathrm{req}} \cdot m}{T},
    -1, 1\right)\right)
\end{equation}

A corrected torque command is then computed using the inner attitude
loop:

\begin{equation}
    \tau_\mathrm{safe} = I\left(k_{p,\theta}(\theta_\mathrm{safe}
    - \theta) + k_{d,\theta}(0 - \dot{\theta})\right)
\end{equation}

Thrust is left unchanged — the CBF only modifies attitude.

\subsection{Verification Test}

\begin{lstlisting}[language=Python]
from params import PARAMS
from cbf import CBFFilter

cbf = CBFFilter(PARAMS)

# Fast approach toward wall: vx = +1.6 m/s
fast_state = [1.85, 1.5, 1.6, 0.0, 0.0, 0.0]
T_safe, tau_safe, active = cbf.filter(fast_state, 9.81, 0.0, PARAMS)
print("CBF active:", active)          # True
print("tau_safe:", round(tau_safe, 4)) # non-zero

# Static hover: no intervention needed
static_state = [1.85, 1.5, 0.0, 0.0, 0.0, 0.0]
T_safe2, tau_safe2, active2 = cbf.filter(static_state, 9.81, 0.0,
                                          PARAMS)
print("CBF active:", active2)         # False
print("T_safe:", round(T_safe2, 4))   # 9.81
\end{lstlisting}

\textbf{Result:}
\begin{itemize}
    \item Fast approach: \texttt{CBF active: True},
          \texttt{tau\_safe = -0.0178} (corrective torque applied)
    \item Static hover: \texttt{CBF active: False},
          \texttt{T\_safe = 9.81} (nominal input passed through)
\end{itemize}

% -------------------------------------------------------
\section{Integration into run\_sim.py}
% -------------------------------------------------------

The \texttt{ode()} inner function was updated to follow this sequence
at every timestep:

\begin{enumerate}
    \item Get nominal control $(T_\mathrm{nom}, \tau_\mathrm{nom})$
          from the Phase 1 controller
    \item Compute $M$ and sub-margins
    \item Update state machine to get current mode
    \item If abort triggered, redirect target to
          $x_\mathrm{retreat}$ and recompute nominal control
    \item Apply CBF filter if mode is \texttt{surface},
          \texttt{docking}, or \texttt{abort}
    \item Log all quantities including mode and CBF flag
    \item Return dynamics derivatives using safe control input
\end{enumerate}

A \texttt{SafeCtrl} helper class was used to pass the safe control
input into the existing \texttt{derivatives()} function without
modifying its interface.

% -------------------------------------------------------
\section{Experiments and Results}
% -------------------------------------------------------

\subsection{Experiment 5 --- CBF Activation Test}

\textbf{Purpose:} Verify the CBF filter activates at the correct $M$
threshold and that the full approach-abort-retreat cycle works.

\textbf{Setup:} Approach speed 0.05\,m/s from $x = 1.5$\,m,
\texttt{cbf\_on=True}, duration 35\,s.

\textbf{Result: PASS.}
\begin{itemize}
    \item All four mode transitions observed: \texttt{free} $\to$
          \texttt{surface} $\to$ \texttt{docking} $\to$
          \texttt{abort} $\to$ \texttt{free}
    \item Drone never crossed $d_{\min}$ ($x < 1.95$\,m throughout)
    \item Abort triggered at $t \approx 20$\,s, drone retreated to
          $x \approx 1.4$\,m
    \item $M$ recovered above $M_\mathrm{safe}$ after retreat
    \item Two complete approach-abort-retreat cycles observed in
          35 seconds
\end{itemize}

\subsection{Experiment 6 --- State Machine Transitions}

\textbf{Purpose:} Verify all four mode transitions occur correctly
under a faster, more aggressive approach.

\textbf{Setup:} Approach speed 0.10\,m/s from $x = 1.5$\,m,
\texttt{cbf\_on=True}, duration 30\,s.

\textbf{Result: PASS.}
\begin{itemize}
    \item Complete mode sequence observed twice within 30 seconds
    \item Each transition occurred at the correct $M$ threshold
    \item Drone never crossed $d_{\min}$ despite faster approach
    \item Two abort retreats visible in trajectory at $t \approx
          15$\,s and $t \approx 21$\,s
    \item $M$ recovered above $M_\mathrm{safe}$ after each retreat
\end{itemize}

\subsection{Experiment 7 --- CBF vs.\ No-CBF Comparison}

\textbf{Purpose:} Quantify the safety benefit of the CBF layer by
comparing identical approach scenarios with and without CBF.

\textbf{Setup:} Two runs with identical initial conditions
(approach speed 0.10\,m/s, $x_0 = 1.5$\,m, duration 25\,s):

\begin{center}
\begin{tabular}{lll}
\toprule
\textbf{Run} & \textbf{cbf\_on} & \textbf{File} \\
\midrule
No CBF & False & \texttt{exp7\_no\_cbf} \\
CBF    & True  & \texttt{exp7\_cbf} \\
\bottomrule
\end{tabular}
\end{center}

\textbf{Result: PASS.}
\begin{itemize}
    \item Both trajectories identical until $t \approx 10$\,s when
          CBF begins activating
    \item CBF run: abort triggered at $t \approx 15$\,s, drone
          retreated safely and re-approached
    \item No-CBF run: drone remained dangerously close to wall
          with $M$ touching $M_\mathrm{abort}$ at $t \approx 21$\,s
          and $t \approx 25$\,s with no recovery
    \item $M$ curves diverge clearly after $t \approx 10$\,s,
          confirming the CBF is actively modifying behavior
\end{itemize}

% -------------------------------------------------------
\section{Challenges Encountered}
% -------------------------------------------------------

\subsection{CBF Zero Gradient at Level Flight}

The first CBF implementation used a \texttt{cvxpy} QP with thrust
and torque as decision variables. The wall clearance constraint was:

\begin{equation}
    \left[\frac{\sin\theta}{m},\; 0\right] u \geq \text{RHS}
\end{equation}

When $\theta = 0$, $\sin\theta = 0$, making the left-hand side
identically zero regardless of $u$. The QP returned the nominal
input unchanged in all cases. The fix was to abandon the QP approach
and instead implement the CBF directly as a pitch angle limiter,
which operates at the correct level of the cascaded controller
architecture.

\subsection{Drone Passing Through Wall}

Early runs showed the drone's $x$ coordinate exceeding $x_w = 2.0$\,m
during abort maneuvers. This occurred because the abort retreat
generated a large overshoot that flung the drone past the wall. Two
fixes were applied:

\begin{enumerate}
    \item A wall clamp was added to \texttt{dynamics.py} to prevent
          $x$ from exceeding $x_w - d_{\min}$
    \item The retreat target was changed from $x_\mathrm{retreat}
          = 1.2$\,m to $1.5$\,m to reduce overshoot
\end{enumerate}

\subsection{Abort Not Preventing Re-approach}

After the first abort and retreat, the approach lambda
$x_\mathrm{des}(t) = \min(x_\mathrm{stop},\; x_\mathrm{start}
+ v_\mathrm{approach} \cdot t)$ continued advancing during the
retreat, causing the drone to immediately re-approach the wall once
$M$ recovered. The fix was an \texttt{abort\_triggered} flag that
freezes the target at $x_\mathrm{retreat}$ until $M$ recovers
sufficiently above $M_\mathrm{safe}$.

\subsection{T\_safe Dropping to Zero}

Early CBF gain values of $\gamma = 2.0$ caused $T_\mathrm{safe}$ to
drop to zero when the drone was very close to the wall. This would
have caused the drone to fall. The gains were reduced to
$\gamma_d = 0.5$ and $\gamma_v = 0.5$, which produces safe but
non-zero thrust corrections.

\subsection{Positive vs.\ Negative vx Sign Convention}

During CBF verification, the test state used $v_x = -1.6$\,m/s
expecting this to represent approach toward the wall. However the
wall normal velocity is $v_n = -v_x$, so $v_x = -1.6$ actually
means $v_n = +1.6$ (receding). The correct test state for approach
uses $v_x = +1.6$\,m/s, giving $v_n = -1.6$ (closing).

% -------------------------------------------------------
\section{Final Parameter Values}
% -------------------------------------------------------

\begin{center}
\begin{tabular}{lll}
\toprule
\textbf{Parameter} & \textbf{Final Value} & \textbf{Description} \\
\midrule
\texttt{cbf\_gamma\_d}  & 0.5  & Wall clearance class-$\mathcal{K}$ gain \\
\texttt{cbf\_gamma\_v}  & 0.5  & Velocity class-$\mathcal{K}$ gain \\
\texttt{cbf\_gamma\_th} & 1.0  & Attitude class-$\mathcal{K}$ gain \\
\texttt{cbf\_gamma\_u}  & 1.0  & Actuator class-$\mathcal{K}$ gain \\
\texttt{M\_safe}        & 0.6  & CBF activates below this value \\
\texttt{M\_dock}        & 0.4  & Docking mode entry \\
\texttt{M\_abort}       & 0.15 & Abort and retreat triggered \\
\texttt{x\_retreat}     & 1.5  & Retreat target position (m) \\
\texttt{retreat\_speed} & 0.05 & Retreat approach speed (m/s) \\
\bottomrule
\end{tabular}
\end{center}

% -------------------------------------------------------
\section{Phase 2 Success Checklist}
% -------------------------------------------------------

\begin{center}
\begin{tabular}{cl}
\toprule
\textbf{Status} & \textbf{Item} \\
\midrule
$\checkmark$ & \texttt{cvxpy} installs and imports without error \\
$\checkmark$ & \texttt{state\_machine.py} shell test passes all four
               modes \\
$\checkmark$ & \texttt{cbf.py} shell test shows
               \texttt{cbf\_active=True} near wall \\
$\checkmark$ & Experiment 5 runs to completion with no errors \\
$\checkmark$ & All four mode transitions observed in Experiment 5 \\
$\checkmark$ & Drone never crosses $d_{\min}$ with CBF on \\
$\checkmark$ & Abort retreat and $M$ recovery working correctly \\
$\checkmark$ & Experiment 6 shows all four mode transitions under
               faster approach \\
$\checkmark$ & Multiple complete approach-abort-retreat cycles
               observed \\
$\checkmark$ & Experiment 7 shows CBF prevents unsafe wall proximity
               \\
$\checkmark$ & No-CBF run shows $M$ touching $M_\mathrm{abort}$
               without recovery \\
$\checkmark$ & All experiment logs saved as \texttt{.npz} files \\
$\checkmark$ & Mode transition plot saved to \texttt{plots/} \\
$\checkmark$ & CBF vs.\ No-CBF comparison plot saved to
               \texttt{plots/} \\
\bottomrule
\end{tabular}
\end{center}

% -------------------------------------------------------
\section{Notes for Phase 3}
% -------------------------------------------------------

Phase 3 moves from 2D planar simulation to a full 3D model and
begins preparation for hardware integration with the PX4 autopilot
stack. The following Phase 2 outputs feed directly into Phase 3:

\begin{itemize}
    \item The CBF pitch angle limiter architecture is directly
          extensible to 3D by adding roll and yaw barrier functions
          alongside the existing pitch barrier.
    \item The calibrated $M$ thresholds
          ($M_\mathrm{safe} = 0.6$, $M_\mathrm{dock} = 0.4$,
          $M_\mathrm{abort} = 0.15$) carry over unchanged.
    \item The state machine mode structure carries over unchanged.
          Only the transition logic needs extension for 3D attitude.
    \item Experiment 7 data provides the baseline for comparing
          simulation CBF performance against hardware performance
          in Phase 4.
\end{itemize}

\textbf{Key limitations to address in Phase 3:}

\begin{itemize}
    \item The 2D planar model does not capture roll dynamics or
          lateral wall effects. These will become significant in
          hardware testing.
    \item The CBF currently modifies only torque. In 3D, thrust
          vectoring effects will require thrust modification as well.
    \item The z-axis drift observed throughout Phase 2 (z settling
          at 1.58\,m instead of 1.5\,m) is caused by wall-induced
          thrust modification $\Gamma(d)$. A feedforward compensation
          term should be added in Phase 3.
    \item The Gazebo simulation environment and PX4 SITL stack should
          be installed before beginning Phase 3.
\end{itemize}

\end{document}